% !TeX root = ../抽象代数笔记.tex
% 由主文档通过 \input 载入；请编译主文档。
\setcounter{chapter}{11} % 保留原书章号。
\chapter{整除性与唯一分解}\label{ch:factorization}

\begin{description}
  \item[学习目标] 识别整数分解论中哪些性质能够推广到整环，区分不同强度的分解条件。
  \item[先修知识] 第\ref{ch:ideals}章。
\end{description}

\section{整除、相伴、不可约元与素元}
本章设 $R$ 为整环。定义整除关系、相伴元（associates）、不可约元（irreducible element）与素元（prime element），两类特殊元素均要求非零且不是单位。证明素元必不可约，说明一般整环中逆命题不成立。最大公因数若存在，仅在相伴意义下唯一。

\section{欧几里得整环与主理想整环}
定义欧几里得整环（Euclidean domain）及其在非零元素上的取值于非负整数的除法函数，要求每次带余除法的非零余式具有更小函数值。定义主理想整环（principal ideal domain, PID），证明欧几里得整环必为主理想整环，并由理想 $(a,b)$ 推导最大公因数的贝祖表示。

\section{唯一分解整环及其判别}
定义唯一分解整环（unique factorization domain, UFD）：每个非零非单位均能分解为有限个不可约元，分解除顺序与相伴替换外唯一。证明主理想整环必为唯一分解整环，分别处理分解的存在与唯一性；总结“欧几里得整环 $\Rightarrow$ 主理想整环 $\Rightarrow$ 唯一分解整环”，避免未经证明倒用箭头。

\section{分式域与典型反例}
对整环 $R$，用分母非零的有序对及交叉相乘等价关系构造分式域（field of fractions）$\operatorname{Frac}(R)$。以高斯整数中的范数说明欧几里得算法，以 $\mathbb Z[\sqrt{-5}]$ 中 $6=2\cdot3=(1+\sqrt{-5})(1-\sqrt{-5})$ 分析唯一分解失效；必须验证各因子不可约且不相伴，而不只展示两种乘积。

\section{例题、练习与自检}
在 $\mathbb Z[i]$ 中做带余除法并求最大公因数，练习从生成理想转化为整除关系。将 $F[x,y]$（$F$ 为域）作为唯一分解整环不必为主理想整环的例子，其唯一分解性待第\ref{ch:polynomials}章建立多项式分解定理后补证；本章先分析理想 $(x,y)$ 为什么不是主理想。

